%\documentclass[a4paper,12pt]{article}
\documentclass{amsart}
\usepackage{amsthm,amsmath,amssymb,amscd,epic}%german,psfrag,

\usepackage[mathscr]{eucal}
 

\newcommand{\cA}{\mathcal{A}}
\newcommand{\cB}{\mathcal{B}}
\newcommand{\cC}{\mathcal{C}}
\newcommand{\cD}{\mathcal{D}}
\newcommand{\cE}{\mathcal{E}}
\newcommand{\cF}{\mathcal{F}}
\newcommand{\cG}{\mathcal{G}}
\newcommand{\cH}{\mathcal{H}}
\newcommand{\cI}{\mathcal{I}}
\newcommand{\cJ}{\mathcal{J}}
\newcommand{\cK}{\mathcal{K}}
\newcommand{\cL}{\mathcal{L}}
\newcommand{\cN}{\mathcal{N}}
\newcommand{\co}{\mathcal O}
\newcommand{\cO}{\mathcal{O}}
\newcommand{\cP}{\mathcal{P}}
\newcommand{\cV}{\mathcal{V}}
\newcommand{\cW}{\mathcal{W}}

\newcommand{\A}{\mathbb A}
\newcommand{\C}{\mathbb C}
\newcommand{\F}{\mathbb F}
\newcommand{\bG}{\mathbb G}
\newcommand{\N}{\mathbb N}
\newcommand{\bP}{\mathbb{P}}
\newcommand{\Q}{\mathbb Q}
\newcommand{\R}{\mathbb R}
\newcommand{\Z}{\mathbb Z}

\newcommand{\kk}{\mathit{k}}
\newcommand{\bs}{\mathbf{s}}
\newcommand{\br}{\mathbf{r}}

\newcommand{\fm}{\mathfrak{m}}
\newcommand{\fp}{\mathfrak{p}}
\newcommand{\fq}{\mathfrak{q}}
\newcommand{\fA}{\mathfrak{A}}
\newcommand{\fB}{\mathfrak{B}}
\newcommand{\fI}{\mathfrak I}
\newcommand{\fJ}{\mathfrak J}
\newcommand{\fR}{\mathfrak{R}} %%%%%%%%%%%%%%%%%%%%%%%%%%%%%%%%%% change: {\mathfrak{P}}
\newcommand{\fT}{\mathfrak T}

\newcommand{\vB}{\check{B}}
\newcommand{\vC}{\check{C}}
\newcommand{\vH}{\check{H}}
\newcommand{\vL}{\check{L}}
\newcommand{\vZ}{\check{Z}}
\newcommand{\vcC}{\check{\mathcal{C}}}

%%%%%%%%%%%%%%%%%%Greek letters%%%%%%%%%%%%%%%%%%%%%%%%%%%%%
\newcommand{\al}{\alpha}
\newcommand{\veps}{\varepsilon}
\newcommand{\Om}{\Omega}
\newcommand{\om}{\omega}
\newcommand{\bmu}{\text{\boldmath{$\mu$}}}
\newcommand{\z}{\zeta}
\newcommand{\G}{\Gamma}
\newcommand{\eps}{\epsilon}
\newcommand{\lam}{\lambda}
\newcommand{\ga}{\gamma}
%%%%%%%%%%%%%%%%%%%%%%%%%%%%%%%%%%%%%%%%%%%%%%%%%%%%%%%%%%%

\newcommand{\Mod}{\mathit{Mod}}
\newcommand{\Ab}{\mathit{Ab}}
\newcommand{\Qcoh}{\mathit{Qcoh}}
\newcommand{\Coh}{\mathit{Coh}}
\newcommand{\coh}{\mathit{coh}}
\newcommand{\rmod}{\mathrm{Mod}}
\newcommand{\rRep}{\mathrm{Rep}}
\newcommand{\rIrr}{\mathrm{Irr}}
\newcommand{\rinj}{\mathrm{inj}}
\newcommand{\rInj}{\mathrm{Inj}}
\newcommand{\rsoc}{\mathrm{soc}}
\newcommand{\JH}{\mathrm{JH}}
\newcommand{\Stab}{\mathrm{Stab}}

\newcommand{\Max}{\text{Max}}


\def\dim{\mathrm{dim}}
\def\Ker{\mathrm{Ker}}
\def\deg{\mathrm{deg}}
\def\Spec{\mathrm{Spec}}
\def\Supp{\mathrm{Supp}}
\def\Ann{\mathrm{Ann}}
\def\Ass{\mathrm{Ass}}
\def\q{\mathfrak {q}}
\def\p{\mathfrak {p}}
\def\m{\mathfrak {m}}
\def\a{\mathfrak {a}}
\def\I{\mathfrak {I}}
\def\x{\chi}
\def\gr{\mathrm{gr}}
\def\pp{\mathbb{P}}
\def\oo{\mathcal{O}}
\def\Ht{\mathrm{ht}} % can't define $\ht$
\def\Fr{\mathrm{Fr}}
\def\St{\mathrm{St}} % the Steinberg representation
\def\Sp{\mathrm{Sp}}
\def\val{\mathrm{val}}
\def\red{\mathrm{red}}
\def\Fil{\mathrm{Fil}}
\def\rg{\mathrm{rg}}

\def\lra{\longrightarrow}
\def\llra{\Longleftrightarrow}
\def\ho{\hookrightarrow}
\def\tw{\twoheadrightarrow}
\def\da{\dashrightarrow}
\newcommand{\calF}{\mathcal{F}}
\DeclareMathOperator{\im}{\mathrm Im}
\DeclareMathOperator{\coker}{\mathrm Coker }
\DeclareMathOperator{\CoIm}{\mathrm CoIm}

\DeclareMathOperator{\End}{\mathrm End}
\DeclareMathOperator{\Hom}{\mathrm Hom}
\DeclareMathOperator{\cHom}{\mathcal{H}\mathit{om}}
\DeclareMathOperator{\cRHom}{\mathrm{R}\cHom}
\DeclareMathOperator{\RHom}{\mathrm {RHom}}
\DeclareMathOperator{\Ext}{\mathrm Ext}
\DeclareMathOperator{\cExt}{\mathcal{E}\mathit{xt}}
\newcommand{\Exti}{\Ext^i}
\newcommand{\cExti}{\cExt^i}
\DeclareMathOperator{\Tor}{\mathrm Tor}
\DeclareMathOperator{\cTor}{\mathcal{T}\mathit{or}}
\DeclareMathOperator{\cEnd}{\mathcal{E}\mathit{nd}}

\DeclareMathOperator{\Frac}{\mathrm Frac}
\DeclareMathOperator{\GL}{\mathrm GL}
\DeclareMathOperator{\SL}{\mathrm SL}
\DeclareMathOperator{\Aut}{\mathrm Aut}
\DeclareMathOperator{\ch}{\mathrm char}
\DeclareMathOperator{\cone}{\mathrm Cone}
\DeclareMathOperator{\depth}{\mathrm depth}
\DeclareMathOperator{\Der}{\mathrm Der}
\DeclareMathOperator{\Et}{\mathrm{Et}}
\DeclareMathOperator{\Gal}{\mathrm Gal}
\DeclareMathOperator{\id}{\mathrm Id}
\DeclareMathOperator{\ord}{\mathrm ord}
\DeclareMathOperator{\Proj}{\mathrm {Proj}}
\DeclareMathOperator{\bProj}{\mathbf{Proj}}
\DeclareMathOperator{\projdim}{\mathrm {proj\,dim}}
\DeclareMathOperator{\qis}{\mathrm qis}
\DeclareMathOperator{\rad}{\mathrm {rad}}
\DeclareMathOperator{\rank}{\mathrm rank}
\DeclareMathOperator{\Sel}{\mathrm Sel }
\DeclareMathOperator{\Spf}{\mathrm{Spf}}
\DeclareMathOperator{\Sym}{\mathrm{Sym}}
\DeclareMathOperator{\tr}{\mathrm Tr}
\DeclareMathOperator{\Rep}{\underline{\mathrm Rep}}
\DeclareMathOperator{\Ind}{\mathrm {Ind}}
\DeclareMathOperator{\ind}{\mathrm {ind}}
\DeclareMathOperator{\cInd}{\mathrm{c-Ind}}
\DeclareMathOperator{\triv}{\mathrm{triv}}
\DeclareMathOperator{\Alg}{\mathrm{Alg}}
\DeclareMathOperator{\Irr}{\mathrm{Irr}}
\DeclareMathOperator{\plim}{\varprojlim}
\DeclareMathOperator{\ilim}{\varinjlim}
%%%%%%%%%%%%%%%%%%%%%%%%%%%%%%%%%% 

\newcommand{\unD}[2][B]{\underline{D}_{#1}({#2})}
\newcommand{\unRep}[2][F]{\underline{\mathrm{Rep}}_{#1}^{#2}(G)}
\newcommand{\fixed}[1][V]{(B\otimes_{F}{#1})^{G}}
\newcommand{\unM}[1][V]{\underline{M}(#1)}
\newcommand{\unV}[1][M]{\underline{V}({#1})}
\newcommand{\BFV}{B\otimes_{F}V}
\newcommand{\BEV}{B\otimes_{E}V}
\newcommand{\xto}[1][]{\xrightarrow{#1}}
\newcommand{\vvec}[2]{\ensuremath{#1_1,\ldots,#1_{#2}}}
\newcommand{\simto}{%\stackrel{\sim}{\to}}%isomorphic to
\xto[\sim]} %isomorphic to;
\newcommand{\EFpV}[1][V]{E^s\otimes_{F_p}{#1}}
\newcommand{\ba}[1][K]{\bar{#1}}
\newcommand{\Kinf}{K_{\infty}}
\newcommand{\Winf}{W_{\infty}}
\newcommand{\cKinf}{\hat{K}_{\infty}}
\newcommand{\cWinf}{\hat{W}_{\infty}}
\newcommand{\Mn}{\mathrm{M}_n}
\newcommand{\cM}{\mathcal{M}}
\DeclareMathOperator{\etale}{\acute{e}tale}
\DeclareMathOperator{\et}{\acute{e}t}
\DeclareMathOperator{\Eur}{\widehat{\cE^{ur}}}
\newcommand{\plimm}{\varprojlim\limits}
\newcommand{\ilimm}{\varinjlim\limits}
\newcommand{\summ}{\sum\limits}
\newcommand{\limm}{\lim\limits}
\newcommand{\prodd}{\prod\limits}
\newcommand{\bigcupp}{\bigcup\limits}
\newcommand{\bigcapp}{\bigcap\limits}
%%%%%%%%%%%%%%%%%%%%%%%%%%%%%%%%%% 
\newcommand{\bFpt}{\overline{\F}_p^{\times}}
\newcommand{\bFp}{\overline{\F}_p}
\newcommand{\bQp}{\overline{\Q}_p}
\newcommand{\ModM}[2][\gamma]{M_{#1}^{#2}}
\newcommand{\ModL}[2][\gamma]{L_{#1}^{#2}}
\newcommand{\matr}[4]{\begin{pmatrix}{#1}&{#2}\\ {#3}&{#4}\end{pmatrix}}
\newcommand{\smatr}[4]{\bigl(\begin{smallmatrix} {#1}& {#2}\\ {#3}&{#4}\end{smallmatrix}\bigl)}
\newcommand{\bmd}{b\hspace{-2mm}\mod d}
\newcommand{\fmm}{\hspace{5mm} }
\newcommand{\gen}[1]{\langle{#1}\rangle_{\bFp}}
\newcommand{\ModV}[2][\bFp]{V_{{#2},{#1}}}
\newcommand{\TModV}[3]{V_{{#1},\bFp}\otimes\ldots\otimes V_{{#2},\bFp}^{\mathrm{Fr}^{#3}}}
\newcommand{\TModVD}[4]{V_{{#1},\bFp}\otimes\ldots\otimes V_{{#2},\bFp}^{\mathrm{Fr}^{#3}}\otimes(\det)^{#4}}
\newcommand{\COEF}[1][G]{\mathcal{COEF}_{#1}}
\newcommand{\DIAG}{\mathcal{DIAG}}
\newcommand{\MOD}{\mathrm{MOD}}
\newcommand{\WD}{\mathrm{WD}}
\newcommand{\rec}{\mathrm{rec}}
\newcommand{\unit}{\mathrm{unit}}
\newcommand{\diag}{\mathrm{diag}}


\newcommand{\ide}{\mathbbm{1}}


%%%%%%%%%%%%%%%%%%%%%%%%%%%%%%%%%%%%%%%%%%%%%%%%%%%%%%%%
\def\dis{\displaystyle}

\newtheorem{theorem}{Th\eb or\ed me}[section]
\newtheorem{lemma}[theorem]{Lemme}
\newtheorem{cor}[theorem]{Corollaire}
\newtheorem{prop}[theorem]{Proposition}


\theoremstyle{definition}
\newtheorem{defn}[theorem]{D\eb finition}
\newtheorem{conj}[theorem]{Conjecture}
\newtheorem{question}[theorem]{Question}
\newtheorem{rem}[theorem]{Remarque}
\newtheorem{examples}[theorem]{Exemples}

%%%%
\newcommand{\ra}{\rightarrow}
\newcommand{\la}{\leftarrow}
\newcommand{\lla}{\longleftarrow}
 
\newcommand{\psp}{\Psi^+}
\newcommand{\psm}{\Psi^-}
\newcommand{\php}{\Phi^+}
\newcommand{\phm}{\Phi^-}
\newcommand{\phhp}{\hat{\Phi}^+}
\newcommand{\vv}{\vspace{2mm}}


\title[]{ \large {\rm Report}    }
\date{}


\begin{document}

\maketitle
 
 
 

 

The paper under consideration is a contribution to the mod $p$ Langlands program for $\GL_2(L)$, where $L$ is a finite unramified extension of $\Q_p$. Currently, the correspondence is only well-understood for $\GL_2(\Q_p)$. One promising approach to the case $L\neq \Q_p$  is via the  local-global compatibility philosophy which, for  a given residual local representation $\overline{\rho}$ and its globalization $\overline{r}$,  provides a candidate $\GL_2(L)$-representation $\pi_{\rm glob}(\overline{\rho})$ in  spaces of mod $p$   automorphic forms.  However, one of the drawbacks of this construction is that $\pi_{\rm glob}(\overline{\rho})$ could depend on the globalization $\overline{r}$ and also various global choices.
\vv

In the paper, the authors show that, at least, a small piece of $\pi_{\rm glob}(\overline{\rho})$   indeed depends only on $\overline{\rho}$, namely the $0$-diagram $\mathcal{D}(\overline{\rho}):=(\pi_{\rm glob}(\overline{\rho})^{I_1}\hookrightarrow \pi_{\rm glob}(\overline{\rho})^{K_1})$ in the sense of Breuil-Pa\v{s}k\=unas. Note that the underlying space $\pi_{\rm glob}(\overline{\rho})^{K_1}$ depends only on $\overline{\rho}$ by a previous result of one of the authors (D.L.) (based on previous work of Emerton-Gee-Savitt). 
The determination of $\mathcal{D}(\overline{\rho})$ is equivalent to determining the values of a certain family of parameters. For the proof of the above independence, they don't (need) compute them explicitly, but rather  identify the values with elements in suitable Galois deformation rings of $\overline{\rho}$.   When  $\overline{\rho}$ is semi-simple, the authors compute explicitly  the values of (some of) the parameters, which    
answers a question of Breuil [J. Reine Angew. Math., 2014]. This result is important, as is related to the structure of $(\varphi,\Gamma)$-module attached to $\mathcal{D}(\overline{\rho})$, and will play a role in the  $p$-adic Langlands program in this setting. 
\vv

 The results in the paper are very nice and shall be proved useful in the future. The paper is very clearly written, despite the technical complexity.   \\
  \textbf{Conclusion}: I strongly recommend the paper for acception in Compositio Math. 
  \bigskip
\bigskip

\noindent Here is a list of  remarks/comments.\bigskip

\vv

 \noindent\textbf{Page 1}
 
Line 4 in the introduction: need to define $G_L=\Gal(\overline{L}/L)$
\vv


 \noindent\textbf{Page 5}
 
 Line 1: do you assume $(x_j)_j\neq (p-1)_j$ ?
 
\vv

 \noindent\textbf{Page 7}
 
 Line 8: the summation should be  over $1\leq t\leq k$.
 
 Line -10: in the formula, it is better to write $S_{i_1}^+S_{i_2}^+$ (exchange the order)
 
 Line -9: typo on $v_p(\alpha)$
 
Line -4: in the formula the last index should be $i_{k+1}$ 

\vv
 \noindent\textbf{Page 9}
 
 

Line 5 of Section 2.2: $\tau|_H$ should be $\theta|_H$

\vv
 \noindent\textbf{Page 10}

Line -11: (3) should be (iii)

Line -9: from Lemma 2.8, what I get is $\mathbf{J}_0(i(\xi_2)+i(\xi_1)+r, -i(\xi_1)-r)$. 
 
\vv

 \noindent\textbf{Page 11}

Line 16: $\tau^1,\tau^2$ should be $\theta^1,\theta^2$
 
  
Line -8: (I think) $S_0\varphi^{\chi^s}$ should be $S_0\varphi^{\chi}$
  
\vv
 \noindent\textbf{Page 12}
 
Line 1, 2: $\sigma(\chi^s)$ should be $\sigma(\xi)$; Line 3: $\chi^s$ should be $\xi$ 

The proof of Lemma 2.14: recall the normalization of $|\cdot|$, i.e. the value of $|p|$.

\vv
 \noindent\textbf{Page 13}
 
Line 2, Line 9, Line 21: $\GL_2$ should be $\GL_n$ 
 
  
\vv 
  \noindent\textbf{Page 15}

Line 4: the citation should be [GHS, Lemma 9.2.4].

Line 10: for $\Lambda_W$, maybe you mean $X(\mathbf{G}_{0}^{\rm der})$ ?

\vv

 \noindent\textbf{Page 21}
 
 Line -7:  what is $\omega_J'$ here ?
 

 \vv
 \noindent\textbf{Page 26}
 
 Line 12: if and only ``if''  

\vv

 \noindent\textbf{Page 29}
 
Line -4: $f_j(\tau)$ should be $f_j(J)$
 
 
\vv

 \noindent\textbf{Page 32}
 
Prop. 4.4: Give a reference for this result. 

After Prop. 4.4: Could you give an example about the possible cycles, say for $f=2$ or $3$ ?
  

\vv
 \noindent\textbf{Page 33}
 
Line 15: Let $\overline{\rho}...$ ``be''  
\vv

%%%%%%
 \noindent\textbf{Page 34}

Line 3:  the tensor product should be taken over  $R_{\infty}(\tau,\lambda)$ 

Line 8: should $K$ be $\GL_2(\cO_L)$? In \S1.1, you defined $K$ to be $\GL_2(\cO_L)$; however, $K$ also denotes the field $L(\pi)$ (\S3.1). Maybe that is why $\GL_2(\cO_L)$ is used in most cases. So, either you make clear this possible confusion of notation, or you use throughout $\GL_2(\cO_L)$. 

\vv
 \noindent\textbf{Page 35}

Lemma 4.9: $i(\chi^s)$ is only used in the proof of Lemma 2.10; you should write $i(\chi^s,R\chi^s)$ instead of $i(\chi^s)$.

Lemma 4.10: It is not clear what the morphism $M_{\infty}(\theta(\chi^s))\ra M_{\infty}(\theta(\chi))$ (until reading the proof), because $\Pi$ is an endomorphism of $M_{\infty}$ but not on $\theta(\chi^s)$ or $\theta(\chi)$. So I suggest first explain the construction of this morphism, using $\Pi$ and Frobenius reciprocity, then state and prove Lemma 4.10.

The proof of Lemma 4.10: in the displayed formula, $\GL_2(k_L)$ should be $\GL_2(\cO_L)$ ($2$ times), because $M_{\infty}^{\vee}$ is not a representation of $\GL_2(k_L)$.

\vv
 \noindent\textbf{Page 36}
 
Line -9: it should be $(M_{\infty}[\mathfrak{m}])^{K_1}\cong D_0(\overline{\rho})$

\vv
 \noindent\textbf{Page 37}
 
The line above the last displayed formula: $\sigma(\chi^s)^{\chi}$ should be $\theta(\chi^s)^{\chi}$ 

\vv

 \noindent\textbf{Page 38}
 
Prop. 4.14: I think you need take $Q(\chi^s)^{R\chi^s}$ to be a quotient of $\theta(\chi^s)^{R\chi^s}/\varpi_E$, i.e. $k_E$-representation ?   (otherwise I don't understand the quotient maps of Line -3).

\vv
 \noindent\textbf{Page 39}
 
 Line -3: $\tau(\chi^s)$ should be $\theta(\chi^s)$; the same remark for Page 40, Line 4.

\vv
 \noindent\textbf{Page 41}
 
The line above the big diagram: $\chi_0,..., \chi_{k-1}$ (miss one comma)

Line -8: in the sequence of morphisms, $\theta^{R\chi_i}$ should be $\theta^{R\chi_i^s}$.

\vv
 \noindent\textbf{Page 42}

Lemma 5.2: since $\overline{\rho}$ can be reducible non-split (until Lemma 5.5), you need to take $J_1\subset J_{\overline{\rho}}$.

\vv
\noindent\textbf{Page 43}

Line -1,,-3: in the formulae, $\tau$ should be $\tau(\chi)$; also $\overline{\mathfrak{M}}^{\tau}(\chi)$ should be $\overline{\mathfrak{M}}^{\tau(\chi)}$

 
\vv

\noindent\textbf{Page 44}

In the middle, you label the assumptions as (i), (ii), (iii); but in the proof of Lemma 5.6, (iii) is cited as (3). 

\vv
\noindent\textbf{Page 46}

Line 5: forgot to precise which $\S$

In the statement of Lemma 5.11, replace $W(k_L)\backslash 0$ by $W(k_L)\backslash\{0\}$

 \vv   
 \noindent\textbf{Page 47}
 
In Proposition 5.14, it is better to cite (32), rather than (23,24).
    \vv
 
 \noindent\textbf{Page 48}

Line 16: in the displayed formula, $\psi_k$ should be $\psi_{k-1}$ (even you identify $\psi_0$ with $\psi_k$)

Section 5.3: In the formula  of $\overline{\rho}$, $\omega_f$ should be $\omega_{2f}$   

\vv
 \noindent\textbf{Page 49}
 
 
Prop. 5.19: $i(\chi_i)^+$ should be $i^+(\chi_i)$ 

Remark 5.20: the citation from [Hu16] should be Lemme 5.4(ii)

Lemma 5.21: it is better to recall what $\mu(t)$ is for $t\in[1,p-2]$.
    
\vv
 \noindent\textbf{Page 50}

Line 12-13: the sentence is broken to two lines.

Line -6: $e(\mu_i)$ should be $e(\mu)_i$ 

\vv
 \noindent\textbf{Page 51}
 
Lemma 5.22: I suggest to write $\mu\lambda$   as $\mu\circ\lambda$, to be compatible with [BP12].

Line -7 (typo): \textbf{P}roposition;   delete ``is''

\vv 
 \noindent\textbf{Page 52}
 
Line 4, 8: replace $1\leq t\leq k$ by $0\leq t\leq k-1$ 
   
\vv

 \noindent\textbf{Page 53}
 
Line 11: $=$ should be $\neq$     

\vv
 \noindent\textbf{Page 55}
 
 In (56), maybe you don't want to put $\prod_{j=0}^{f-1}$ ? 
 
\vv
 \noindent\textbf{Page 56}
 
 Line 11-12: you cite the proof of [Bre11. Proposition 5.1], but I think it is better to also mention [BP12, Lemma 15.2].
 
Line -4: I think you want to cite equation (52) (not Proposition)   
 
\vv
 \noindent\textbf{Page 57}
 
 Line 7 (typo): ``an'' should be ``and''
 
\vv
 \noindent\textbf{Page 59}
 
The end of proof of Thm. 5.36: exchange the order of reducible and irreducible case, to be compatible. Last line: replace $\mu_1-\mu_2$ by $\mu_{1,i}-\mu_{2,i}$. 
 
Section 6.1, first sentence: ``Llet" should be ``Let'' 
 
  
\vv \noindent\textbf{Page 66}

Line 1: ``serre'' should be ``Serre''  
 
 \end{document} 